%-----------------------------------------------------------------------
\section{Introduction}
\label{sec:introduction}
%-----------------------------------------------------------------------

The foundational pillars of modern digital security---specifically RSA and Elliptic Curve Cryptography (ECC)---are mathematically fragile against cryptographically relevant quantum computers (CRQCs). Shor's algorithm \cite{shor1994algorithms} proves that asymmetric cryptographic primitives relying on integer factorization and discrete logarithm problems can be broken in polynomial time. While the engineering timeline for a large-scale, fault-tolerant CRQC remains probabilistic, the global threat is already active due to the ``harvest now, decrypt later'' attack paradigm. In this model, sophisticated adversaries are currently intercepting and storing classical TLS traffic, awaiting quantum maturation to decrypt sensitive, long-term confidential data retroactively \cite{mosca2018cybersecurity}. 

In response to this existential-level cryptographic threat, the National Institute of Standards and Technology (NIST) concluded its multi-year standardization program on August 13, 2024, which finalized the initial set of critical post-quantum cryptography (PQC) standards: FIPS 203 (ML-KEM), FIPS 204 (ML-DSA), and FIPS 205 (SLH-DSA) \cite{nist2024pqc}. This marks the beginning of the global transition to quantum-resistant cryptography.

The Internet of Things (IoT), in particular, is a challenging and vulnerable deployment platform for the transition to post-quantum cryptography. It is forecasted that by 2025, the IoT domain will comprise 75 billion devices, which will form the IoT ecosystem \cite{statista2021iot}. These devices will be responsible for managing critical infrastructures such as medical telemetry, industrial control systems, automotive V2X networking, and smart grids. IoT devices, unlike cloud servers and mobile devices, operate in environments defined by RFC 7228 \cite{bormann2014coap}. These environments are characterized by limited memory available for the cryptographic stack, i.e., lower than 50 KB, microcontrollers running at lower than 50 MHz, and battery lifecycles spanning several years, implying power budgets lower than one milliwatt.

\begin{figure}[htbp]
\centering
\includegraphics[width=\columnwidth]{fig15_migration_cost}
\caption{The severe baseline computational and networking payload expansion characterizing the post-quantum transition.}
\label{fig:migration_cost}
\end{figure}

However, the transition from classical cryptography to post-quantum cryptography, which is based on lattices, hash trees, and codes, is associated with considerable overhead, which is in direct conflict with IoT devices. This is because public keys and signatures in post-quantum cryptography are significantly larger than ECC-based public keys and signatures, resulting in fragmentation of the TLS handshake at the network stack. Moreover, the complex matrix calculations in post-quantum cryptography accelerate the depletion of device power.

\subsection{Novelty and Research Gap}

In spite of the existing foundational literature detailing the performance of individual PQC schemes in terms of embedded systems, such as the pqm4 project's study of Cortex-M4 microcontrollers \cite{kannwischer2019pqm4}, Fitzgibbon and Ottaviani's exhaustive IoT benchmarking study for IoT applications \cite{fitzgibbon2024benchmarking}, and real-world deployment assessments for ESP32 microcontrollers \cite{secrypt2025esp32}, a critical synthesis gap still remains. This is particularly true in terms of a unified, cross-class, systematic evaluation and mapping of the entire gamut of primary PQC schemes, i.e., \kyber{}, \dilithium{}, \falcon{}, \sphincs{}, and \ntru{}, to the exact definitions of the hardware classes defined in RFC 7228 Class 0, Class 1, and Class 2 devices.

In addition, existing literature is heavily weighted towards execution time while largely failing to assess unified models of energy consumption, i.e., computation and radio transmission, and deterministic multi-criteria algorithm selection guidelines.

In terms of addressing the above unmapped intersections, the following questions are examined in this paper: With the latest standardized algorithms defined by NIST, what schemes are viable for specific classes of IoT devices, and what is the relative cost to classical cryptography, and what are the trade-offs?

\subsection{Contributions}

This paper address this gaps and our contributions are:
\begin{enumerate}
 \item \textbf{Comprehensive Cross-Class Benchmarking:} We perform and evaluate more than 45,000 empirical measurements across 14 algorithm parameter classes and correlate them with tightly constrained C0, C1, and C2 hardware simulations.
 \item \textbf{Expanded Energy and Communication Analysis:} We propose a unified energy estimation model that accounts for microcontroller cycle drain and radio frequency transmission overheads of inflated PQC key sizes.
 \item \textbf{TOPSIS Decision Framework:} We apply a multi-attribute statistical decision matrix to objectively evaluate and classify algorithms in terms of speed, memory footprint, energy, key size, and security scale to yield the first deterministic PQC algorithm deployment guide targeting IoT system architects.
 \item \textbf{Classical vs. Quantum and Hybrid Profiling:} Classical bounds are established using RSA and ECC algorithms to determine migration costs, including overhead metrics of Hybrid TLS 1.3 models to bridge the classical-quantum divide.
\end{enumerate}
